\documentclass[11pt]{article}
\usepackage{amsmath,amssymb,amsthm}
\usepackage{lmodern}
\usepackage{geometry}
\usepackage{booktabs}
\usepackage{array}
\usepackage{hyperref}
\usepackage{bm}
\usepackage{bbm}
\usepackage[parfill]{parskip}

\geometry{margin=1.2in}

\newtheorem{theorem}{Theorem}
\newtheorem{proposition}[theorem]{Proposition}
\newtheorem{corollary}[theorem]{Corollary}
\newtheorem{lemma}[theorem]{Lemma}
\theoremstyle{remark}
\newtheorem{remark}{Remark}

\title{%
  \textbf{Matching Estimators and Linear Regression}\\[6pt]
  \large Angrist (1998): Proof for the Saturated and General Linear Cases\\[4pt]
  \normalsize ECON 712 --- Supplementary Notes
}
\author{}
\date{}

\begin{document}
\maketitle
\tableofcontents
\bigskip

%--------------------------------------------------------------------
\section{Setup and Notation}
%--------------------------------------------------------------------

Let $D_i \in \{0,1\}$ be a binary treatment, $X_i$ a vector of pre-treatment covariates,
and $\{Y_{0i}, Y_{1i}\}$ the potential outcomes.  The observed outcome is
$Y_i = Y_{0i} + (Y_{1i} - Y_{0i})D_i$.

Define the following population objects:
\begin{align*}
  p(x)      &\;=\; E[D_i \mid X_i = x]               &&\text{(propensity score)}\\
  \mu_0(x)  &\;=\; E[Y_{0i} \mid X_i = x]             &&\text{(baseline outcome CEF)}\\
  \delta(x) &\;=\; E[Y_{1i} - Y_{0i} \mid X_i = x]   &&\text{(conditional ATE)}
\end{align*}

The true data-generating process can therefore be written as
\begin{equation}\label{eq:dgp}
  Y_i \;=\; \mu_0(X_i) \;+\; \delta(X_i)\,D_i \;+\; u_i,
  \qquad E[u_i \mid D_i, X_i] = 0,
\end{equation}
where the zero-mean condition on $u_i$ follows from the \textbf{Conditional Independence
Assumption (CIA)}: $(Y_{0i},Y_{1i}) \perp\!\!\!\perp D_i \mid X_i$.

\bigskip

\textbf{What we want to know.}  We run a regression that \emph{linearly controls} for
$X_i$ --- either a fully saturated specification or a plain linear-in-$X$ specification.
In both cases we use the Frisch--Waugh--Lovell (FWL) theorem to isolate the coefficient
on $D_i$.  The question is: what causal parameter does that coefficient identify under CIA?

%--------------------------------------------------------------------
\section{The Frisch--Waugh--Lovell Starting Point}
%--------------------------------------------------------------------

Consider any regression of the form
\[
  Y_i \;=\; \alpha \;+\; \delta\,D_i \;+\; g(X_i) \;+\; \varepsilon_i
\]
for some control function $g(\cdot)$.  By FWL, the population coefficient on $D_i$ is
\begin{equation}\label{eq:fwl}
  \operatorname{plim} \hat\delta
  \;=\; \frac{E[\tilde D_i\, Y_i]}{E[\tilde D_i^{\,2}]},
\end{equation}
where $\tilde D_i = D_i - \hat p(X_i)$ and $\hat p(X_i) = L[D_i \mid g(X_i)]$ is the
linear projection of $D_i$ on the control function.

The two cases we study differ \emph{only} in what $\hat p$ is:
\begin{itemize}
  \item \textbf{Saturated regression}: $g(X_i)$ includes a dummy for every cell of
        $X_i$, so the projection is the conditional expectation:
        $\hat p^{\mathrm{sat}}(X_i) = p(X_i)$ exactly.
  \item \textbf{Linear regression}: $g(X_i) = \bm\beta' X_i$, so
        $\hat p^{\mathrm{lin}}(X_i) = \Pi_0 + \bm\Pi' X_i$ is only the best linear
        approximation to $p(X_i)$; in general $\hat p^{\mathrm{lin}} \neq p$.
\end{itemize}

%--------------------------------------------------------------------
\section{The Saturated Case}
%--------------------------------------------------------------------

This is the proof given (implicitly) by Angrist (1998) and made explicit in
Angrist and Pischke (2009, Ch.\ 3).

When the regression is saturated, $\tilde D_i^{\mathrm{sat}} = D_i - p(X_i)$.

\subsection*{Denominator}

\[
  E\!\left[(\tilde D_i^{\mathrm{sat}})^2\right]
  = E\!\left[(D_i - p(X_i))^2\right]
  \stackrel{\mathrm{LIE}}{=} E\!\left[\operatorname{Var}(D_i \mid X_i)\right].
\]

\subsection*{Numerator: Term A --- baseline outcomes}

\[
  E[\tilde D_i^{\mathrm{sat}}\, Y_{0i}]
  = E\!\left[E\!\left[(D_i - p(X_i))\,Y_{0i} \mid X_i\right]\right].
\]
Under CIA, $Y_{0i} \perp D_i \mid X_i$, so
\[
  E[(D_i - p(X_i))\,Y_{0i} \mid X_i]
  = \underbrace{E[D_i - p(X_i)\mid X_i]}_{=\;0} \cdot E[Y_{0i}\mid X_i] \;=\; 0.
\]

\subsection*{Numerator: Term B --- treatment effect heterogeneity}

Under CIA $E[\delta_i \mid D_i, X_i] = \delta(X_i)$, so the $(\delta_i-\delta(X_i))D_i$
residual satisfies $E[\tilde D_i \cdot (\delta_i - \delta(X_i))D_i \mid D_i,X_i] = 0$.
Thus we only need:
\[
  E[\tilde D_i^{\mathrm{sat}}\,\delta(X_i)\,D_i]
  = E\!\left[\delta(X_i)\,E[(D_i - p(X_i))\,D_i \mid X_i]\right].
\]
For binary $D_i$, $D_i^2 = D_i$, so
\[
  E[(D_i - p(X_i))\,D_i \mid X_i] = p(X_i) - p(X_i)^2
  = p(X_i)(1-p(X_i)) = \operatorname{Var}(D_i\mid X_i).
\]
Hence
\[
  E[\tilde D_i^{\mathrm{sat}}\,\delta(X_i)\,D_i]
  = E\!\left[\delta(X_i)\operatorname{Var}(D_i\mid X_i)\right].
\]

\subsection*{Result}

Assembling \eqref{eq:fwl}:

\begin{proposition}[Angrist 1998 --- saturated regression]\label{prop:sat}
Under CIA and a fully saturated regression,
\[
  \operatorname{plim}\,\hat\delta^{\mathrm{sat}}
  \;=\;
  \frac{E\!\left[\delta(X_i)\operatorname{Var}(D_i\mid X_i)\right]}
       {E\!\left[\operatorname{Var}(D_i\mid X_i)\right]}
  \;=\;
  \int \delta(x)\,\omega(x)\,dx,
\]
where the variance weights are
\[
  \omega(x) \;=\;
  \frac{\operatorname{Var}(D_i\mid X_i=x)\cdot f(x)}
       {E[\operatorname{Var}(D_i\mid X_i)]}, \qquad \int\omega(x)\,dx = 1.
\]
\end{proposition}

\begin{remark}
The weights $\omega(x)$ are largest where the propensity score is closest to $\tfrac12$,
i.e.\ where treatment varies most.  Under CIA, ATT$(x)=$ATU$(x)=$ATE$(x)=\delta(x)$,
so $\hat\delta^{\mathrm{sat}}$ is a weighted average of the ATE across covariate cells.
\end{remark}

%--------------------------------------------------------------------
\section{The General Linear Case}
%--------------------------------------------------------------------

Now suppose we only control linearly: $Y_i = \alpha + \delta D_i + \bm\beta'X_i +
\varepsilon_i$.  The projected propensity score is $\hat p^L(X_i) = \Pi_0 + \bm\Pi'X_i$,
the best linear predictor of $D_i$ given $(1, X_i)$.

\textbf{Key object.}  Define the \emph{propensity-score misspecification residual}
\begin{equation}\label{eq:lambda}
  \lambda(X_i) \;\equiv\; p(X_i) - \hat p^L(X_i).
\end{equation}
This is the component of the propensity score that the linear regression \emph{cannot
capture} --- the nonlinearity in $p(\cdot)$.  The OLS normal equations imply
$E[\lambda(X_i)] = 0$ and $E[\lambda(X_i)\,X_i] = \bm 0$.

The linear regression residual decomposes as
\[
  \tilde D_i^{\mathrm{lin}}
  = D_i - \hat p^L(X_i)
  = \underbrace{(D_i - p(X_i))}_{\displaystyle e_i} + \lambda(X_i),
  \qquad E[e_i\mid X_i]=0.
\]

\subsection{Denominator}

\begin{align*}
  E\!\left[(\tilde D_i^{\mathrm{lin}})^2\right]
  &= E[e_i^2] + 2\,E[e_i\,\lambda(X_i)] + E[\lambda(X_i)^2] \\
  &= E\!\left[\operatorname{Var}(D_i\mid X_i)\right]
     + 2\,E\!\left[\underbrace{E[e_i\mid X_i]}_{=\,0}\lambda(X_i)\right]
     + E[\lambda(X_i)^2] \\[4pt]
  &= E\!\left[\operatorname{Var}(D_i\mid X_i)\right] + \operatorname{Var}(\lambda(X_i)).
\end{align*}
The denominator is \emph{strictly larger} than in the saturated case whenever
$\lambda \not\equiv 0$.

\subsection{Numerator}

From \eqref{eq:dgp}, $Y_i = \mu_0(X_i) + \delta(X_i)D_i + u_i$.

\medskip\noindent\textbf{Term A: baseline outcomes.}
\begin{align*}
  E[\tilde D_i^{\mathrm{lin}}\,Y_{0i}]
  &= E\!\left[E[(D_i - \hat p^L(X_i))\,Y_{0i}\mid X_i]\right]\\
  &= E\!\left[(E[D_i\mid X_i] - \hat p^L(X_i))\,\mu_0(X_i)\right]
     \quad\text{(CIA: }Y_{0i}\perp D_i\mid X_i\text{)}\\
  &= E[\lambda(X_i)\,\mu_0(X_i)].
\end{align*}
Unlike the saturated case, this is \emph{not} zero unless $\lambda \equiv 0$.

\medskip\noindent\textbf{Term B: treatment effects.}

The residual $(\delta_i-\delta(X_i))D_i$ contributes zero by the same CIA argument
as before.  For the systematic part:
\begin{align*}
  E[\tilde D_i^{\mathrm{lin}}\,\delta(X_i)\,D_i]
  &= E\!\left[\delta(X_i)\,E[(D_i-\hat p^L(X_i))\,D_i\mid X_i]\right] \\
  &= E\!\left[\delta(X_i)\,p(X_i)(1-\hat p^L(X_i))\right].
\end{align*}
Decompose $1-\hat p^L = (1-p)+\lambda$:
\[
  = E\!\left[\delta(X_i)\operatorname{Var}(D_i\mid X_i)\right]
    + E\!\left[\delta(X_i)\,p(X_i)\,\lambda(X_i)\right].
\]

\medskip\noindent\textbf{Combining.}  Adding Terms A and B:
\[
  E[\tilde D_i^{\mathrm{lin}}\,Y_i]
  = E\!\left[\delta(X_i)\operatorname{Var}(D_i\mid X_i)\right]
  + E\!\left[\lambda(X_i)\!\underbrace{\bigl(\mu_0(X_i)+\delta(X_i)p(X_i)\bigr)}_{
    =\;E[Y_i\mid X_i]}\right].
\]

\medskip\noindent\textbf{Simplifying the bias term.}
Decompose $E[Y_i\mid X_i] = L[E[Y_i\mid X_i]\mid 1,X_i] + \rho(X_i)$, where
$\rho(X_i)$ is the \emph{nonlinear residual of the outcome CEF}.  Since $\lambda(X_i)$
is orthogonal to all linear functions of $(1,X_i)$ (by the normal equations),
\[
  E[\lambda(X_i)\,E[Y_i\mid X_i]]
  = \underbrace{E[\lambda(X_i)\,L[E[Y_i\mid X_i]\mid 1,X_i]]}_{=\;0}
  + E[\lambda(X_i)\,\rho(X_i)]
  = \operatorname{Cov}(\lambda(X_i),\,\rho(X_i)).
\]

\subsection{Result}

\begin{proposition}[Linear regression --- general case]\label{prop:lin}
Under CIA and the linear regression $Y_i = \alpha + \delta D_i + \bm\beta'X_i +
\varepsilon_i$,
\begin{equation}\label{eq:main}
  \operatorname{plim}\,\hat\delta^{\mathrm{lin}}
  \;=\;
  \frac{%
    E\!\left[\delta(X_i)\operatorname{Var}(D_i\mid X_i)\right]
    \;+\;
    \operatorname{Cov}\!\bigl(\lambda(X_i),\,\rho(X_i)\bigr)
  }{%
    E\!\left[\operatorname{Var}(D_i\mid X_i)\right]
    \;+\;
    \operatorname{Var}(\lambda(X_i))
  },
\end{equation}
where
\begin{align*}
  \lambda(X_i) &\;=\; p(X_i) - L[p(X_i)\mid 1,X_i]
    &&\text{(nonlinear part of propensity score),}\\
  \rho(X_i)   &\;=\; E[Y_i\mid X_i] - L[E[Y_i\mid X_i]\mid 1,X_i]
    &&\text{(nonlinear part of outcome CEF).}
\end{align*}
\end{proposition}

%--------------------------------------------------------------------
\section{Comparing the Two Cases}
%--------------------------------------------------------------------

\subsection{Structural comparison}

\begin{center}
\begin{tabular}{lcc}
  \toprule
  & \textbf{Saturated} & \textbf{Linear} \\
  \midrule
  $\hat p(X_i)$ & $p(X_i)$ & $L[p(X_i)\mid 1,X_i]$ \\
  $\lambda(X_i)$ & $0$ & $p(X_i)-\hat p^L(X_i)$ \\[3pt]
  Denominator & $E[\operatorname{Var}(D_i\mid X_i)]$
              & $E[\operatorname{Var}(D_i\mid X_i)] + \operatorname{Var}(\lambda)$ \\[3pt]
  Numerator bias & $0$
                 & $\operatorname{Cov}(\lambda(X_i),\rho(X_i))$ \\[3pt]
  Weights sum to 1? & Yes & Not in general \\
  \bottomrule
\end{tabular}
\end{center}

\subsection{When does the linear case recover the variance-weighted formula?}

Setting $\lambda \equiv 0$ in \eqref{eq:main} immediately yields Proposition~\ref{prop:sat}.
Three sufficient conditions for $\lambda \equiv 0$ (or for the bias terms to vanish) are:

\begin{enumerate}
  \item \textbf{The propensity score is linear in $X_i$.}  Then
        $\hat p^L(X_i) = p(X_i)$ exactly and $\lambda \equiv 0$.
        The linear regression gives the same result as the saturated regression
        \emph{regardless} of whether $\mu_0(X_i)$ is linear.

  \item \textbf{The outcome CEF is linear in $X_i$.}  Then $\rho \equiv 0$, so
        $\operatorname{Cov}(\lambda,\rho)=0$.  \emph{But} the denominator is still
        inflated by $\operatorname{Var}(\lambda)\geq 0$ unless $\lambda\equiv 0$ as well.
        The pure linear-outcome case does \emph{not} by itself recover the
        variance-weighted formula.

  \item \textbf{Constant treatment effects and linear outcome.}  If
        $\delta(x)=\delta$ for all $x$ and $E[Y_i\mid X_i]$ is linear, a
        calculation shows that the numerator bias equals
        $\delta\,\operatorname{Var}(\lambda)$, so the denominator inflation and
        numerator bias exactly cancel:
        \[
          \operatorname{plim}\,\hat\delta^{\mathrm{lin}}
          = \frac{\delta\,E[\operatorname{Var}(D_i\mid X_i)] + \delta\,\operatorname{Var}(\lambda)}
                 {E[\operatorname{Var}(D_i\mid X_i)] + \operatorname{Var}(\lambda)}
          = \delta.
        \]
        This is why a (correctly specified) linear regression with linear controls is
        consistent for the constant ATE even when the propensity score is nonlinear.
\end{enumerate}

\subsection{Interpreting the bias}

The bias term $\operatorname{Cov}(\lambda(X_i), \rho(X_i))$ has an intuitive reading:

\begin{itemize}
  \item $\lambda(X_i)$ is the part of the propensity score the linear regression
        \emph{fails to absorb} --- cells where the true treatment probability curves
        away from the fitted linear index.
  \item $\rho(X_i)$ is the part of the outcome function the linear regression
        \emph{fails to absorb} --- cells where $E[Y_i\mid X_i]$ is curved.
  \item The bias arises when these two forms of misspecification are
        \emph{correlated}: the cells where the linear regression gets the propensity
        score wrong tend to be the same cells where it also gets the outcome function
        wrong.  OLS then conflates the two, distorting the treatment-effect estimate.
\end{itemize}

\subsection{Why the saturated case avoids the problem}

In the saturated case, $g(X_i)$ contains a dummy for every cell of $X_i$, so the
linear projection of $D_i$ on $g(X_i)$ \emph{is} the cell mean $p(X_i)$.  The
regression is therefore also saturating the propensity-score model.  By construction
$\lambda \equiv 0$, which kills both the numerator bias and the denominator inflation.
There is no residual nonlinearity for $\rho$ to correlate with.

%--------------------------------------------------------------------
\section{Summary}
%--------------------------------------------------------------------

\begin{theorem}[Regression--Matching equivalence]\label{thm:main}
Under CIA:
\begin{enumerate}
  \item[(a)] \textbf{Saturated regression} always identifies the variance-weighted ATE
             of Proposition~\ref{prop:sat}, with weights proportional to the
             conditional variance of treatment.

  \item[(b)] \textbf{Linear regression} identifies the same quantity if and only if
             $\lambda(X_i) \equiv 0$, i.e.\ the true propensity score $p(X_i)$ is
             linear in $X_i$.  Otherwise, the estimand differs by the correction terms
             in Proposition~\ref{prop:lin}.

  \item[(c)] In either case, if treatment effects are constant ($\delta(x)=\delta$)
             \emph{and} the outcome equation is correctly specified, the regression
             identifies the constant ATE $=$ ATT $=$ ATU $= \delta$.
\end{enumerate}
\end{theorem}

\noindent The practical implication is that with heterogeneous treatment effects and
an unknown propensity score, a saturated regression (or a matching estimator using the
nonparametric $p(X_i)$) gives a cleaner and more interpretable weighted average than a
plain linear regression, which may produce weights that are difficult to characterize
and that can even be negative when $p(X_i)$ is strongly nonlinear.

\bigskip
\noindent\textbf{Reference.}
Angrist, J.D.\ (1998). ``Estimating the Labor Market Impact of Voluntary Military
Service Using Social Security Data on Military Applicants.'' \textit{Econometrica},
66(2), 249--288.\\[2pt]
Angrist, J.D.\ and Pischke, J.S.\ (2009). \textit{Mostly Harmless Econometrics}.
Princeton University Press.

\end{document}
